\section{Hạ theo tọa độ và tối ưu luân phiên}
\label{SEC_coordinate_alternation}

\subsection{Ý tưởng luân phiên tối ưu}
\label{SUB_SEC_alternation_idea}
Giả sử ta cần tối thiểu hóa một hàm $f:\mathbb{R}^{n+m}\rightarrow\mathbb{R}$. Thay vì xem biến đầu vào là một vectơ lớn, ta tách nó thành hai khối $x\in\mathbb{R}^{n}$ và $y\in\mathbb{R}^{m}$, rồi viết hàm mục tiêu dưới dạng $f(x,y)$. Một chiến lược đơn giản là cố định $y$ để tối ưu $x$, sau đó cố định $x$ để tối ưu $y$, và lặp lại:
\begin{equation}
    \begin{cases}
        x_{k+1}=\arg\min_x f(x,y_k),\\
        y_{k+1}=\arg\min_y f(x_{k+1},y).
    \end{cases}
    \label{eq_alternation_two_blocks}
\end{equation}
Chiến lược này được gọi là \textbf{tối ưu luân phiên}.
\medskip

Nếu mỗi bài toán con được giải chính xác, giá trị hàm mục tiêu không tăng qua từng bước:
\begin{equation}
    f(x_{k+1},y_k)\leq f(x_k,y_k),
    \qquad
    f(x_{k+1},y_{k+1})\leq f(x_{k+1},y_k).
    \label{eq_alternation_monotone}
\end{equation}
Suy ra
\begin{equation}
    f(x_{k+1},y_{k+1})\leq f(x_k,y_k).
    \label{eq_alternation_energy_decrease}
\end{equation}
Tính giảm đơn điệu này là một ưu điểm lớn của tối ưu luân phiên. Tuy nhiên, điều đó không đảm bảo thuật toán luôn đi tới nghiệm toàn cục. Với bài toán không lồi, tối ưu luân phiên có thể hội tụ tới nghiệm cục bộ hoặc điểm dừng phụ thuộc mạnh vào khởi tạo.

\subsection{Khi nào nên dùng tối ưu luân phiên}
\label{SUB_SEC_when_alternation}
Tối ưu luân phiên thường phù hợp khi bài toán gốc khó, nhưng sau khi cố định một nhóm biến thì bài toán con trở nên đơn giản. Có hai tình huống đặc biệt hay gặp:
\begin{itemize}
    \item Bài toán gốc phi tuyến theo toàn bộ biến, nhưng tuyến tính hoặc có dạng bình phương tối thiểu theo từng khối biến.
    \item Một số biến là biến rời rạc hoặc biến ràng buộc phức tạp, nhưng khi cố định các biến còn lại thì có công thức cập nhật đóng.
\end{itemize}
\medskip

Trong học máy, tối ưu luân phiên xuất hiện rất tự nhiên. Ví dụ, trong k-means, nếu biết tâm cụm thì việc gán nhãn cụm rất dễ; nếu biết nhãn cụm thì việc cập nhật tâm cụm cũng rất dễ. Nhưng nếu tối ưu đồng thời cả tâm cụm và nhãn cụm thì bài toán trở nên khó hơn nhiều.

\subsection{Một số ví dụ tiêu biểu}
\label{SUB_SEC_examples_alternation}

\subsubsection*{PCA tổng quát và bình phương tối thiểu luân phiên}
Cho ma trận dữ liệu $X\in\mathbb{R}^{n\times m}$, trong đó mỗi cột là một điểm dữ liệu. PCA cổ điển tìm một không gian con chiều thấp sao cho dữ liệu được xấp xỉ tốt nhất trong chuẩn Frobenius. Một cách viết là
\begin{equation}
    \min_{C,Y}\|X-CY\|_{\mathrm{Fro}}^2,
    \label{eq_pca_als}
\end{equation}
trong đó $C\in\mathbb{R}^{n\times d}$ chứa các vectơ cơ sở và $Y\in\mathbb{R}^{d\times m}$ chứa hệ số biểu diễn của dữ liệu trong cơ sở đó.
\medskip

Nếu cố định $Y$, bài toán tối ưu $C$ là bình phương tối thiểu. Nếu cố định $C$, bài toán tối ưu $Y$ cũng là bình phương tối thiểu. Vì vậy, theo đúng thứ tự luân phiên trong tài liệu, ta có thuật toán bình phương tối thiểu luân phiên:
\begin{equation}
    \begin{cases}
        C_{k+1}=\arg\min_C \|X-CY_k\|_{\mathrm{Fro}}^2,\\
        Y_{k+1}=\arg\min_Y \|X-C_{k+1}Y\|_{\mathrm{Fro}}^2.
    \end{cases}
    \label{eq_pca_als_steps}
\end{equation}
Nếu thêm ràng buộc trực chuẩn cho các cột của $C$, ta quay về PCA cổ điển. Nếu thay chuẩn Frobenius bằng chuẩn khác, ví dụ chuẩn $L^1$ theo từng phần tử, ta thu được các biến thể robust PCA có khả năng giảm ảnh hưởng của nhiễu và ngoại lai \cite{Candes2011RPCA}.
\medskip

Tích $CY$ làm bài toán không lồi nếu xét đồng thời $C$ và $Y$. Khi một trong hai biến được cố định, phần còn lại giảm về một bài toán đơn giản hơn nhiều.

\subsubsection*{Bài toán ARAP}
ARAP là viết tắt của \textit{as-rigid-as-possible}, thường dùng trong xử lý hình học và biến dạng lưới. Một dạng năng lượng ARAP phẳng có thể viết như sau:
\begin{equation}
    \begin{aligned}
        &\min_{\{R_v\},\{y_v\}}
        \quad
        \sum_{v\in V}\sum_{(v,w)\in E}
        \|R_v(x_v-x_w)-(y_v-y_w)\|_2^2\\
        &\text{với điều kiện}
        \quad
        R_v^{T}R_v=I_2,\quad \forall v\in V,\\
        &\hspace{3.1cm}
        y_v\ \text{cố định},\quad \forall v\in V_0.
    \end{aligned}
    \label{eq_arap_energy}
\end{equation}
Ở đây $x_v$ là vị trí ban đầu, $y_v$ là vị trí sau biến dạng, $R_v$ là ma trận quay cục bộ tại đỉnh $v$, còn $V_0$ là tập các đỉnh được giữ cố định.
\medskip

Nếu tối ưu đồng thời tất cả $R_v$ và $y_v$, bài toán rất khó vì có ràng buộc trực giao. Nhưng nếu tách biến, ta có hai bước:
\begin{itemize}
    \item Cố định các $R_v$, tối ưu các $y_v$. Khi đó bài toán là bình phương tối thiểu thưa, có thể giải bằng hệ tuyến tính xác định dương.
    \item Cố định các $y_v$, tối ưu từng $R_v$. Các bài toán theo từng đỉnh tách rời nhau và trở thành bài toán Procrustes, có thể giải bằng SVD kích thước nhỏ.
\end{itemize}
\medskip

Ví dụ này cho thấy sức mạnh của tối ưu luân phiên: một bài toán phi tuyến lớn được thay bằng một bước giải hệ tuyến tính thưa và nhiều bài toán nhỏ có nghiệm đóng.

\subsubsection*{Hạ theo tọa độ cho bình phương tối thiểu}
Hạ theo tọa độ là trường hợp cực đoan của tối ưu luân phiên, trong đó mỗi lần chỉ tối ưu một tọa độ. Xét bài toán
\begin{equation}
    \min_x \|Ax-b\|_2^2,
    \label{eq_cd_ls}
\end{equation}
với $A=[a_1,a_2,\dots,a_n]$. Khai triển sai số dư:
\begin{equation}
    Ax-b=\sum_{j=1}^{n}x_ja_j-b.
    \label{eq_cd_expansion}
\end{equation}
Khi cố định tất cả các biến trừ $x_i$, điều kiện đạo hàm theo $x_i$ bằng $0$ là
\begin{equation}
    a_i^{T}(Ax-b)=0.
    \label{eq_cd_derivative}
\end{equation}
Suy ra công thức cập nhật
\begin{equation}
    x_i
    \leftarrow
    \frac{a_i^{T}b-\sum_{j\neq i}x_j a_i^{T}a_j}
         {\|a_i\|_2^2}.
    \label{eq_cd_update_ls}
\end{equation}
Thuật toán lặp qua $i=1,2,\dots,n$ nhiều lần cho đến khi hội tụ. Mỗi bước cập nhật chỉ thay đổi một tọa độ, nên chi phí từng bước rất nhỏ. Đặc điểm này đặc biệt hữu ích trong các bài toán có số chiều lớn hoặc có cấu trúc thưa.

\subsubsection*{Thuật toán phân cụm k-means}
Cho các điểm dữ liệu $x_1,x_2,\dots,x_m\in\mathbb{R}^{n}$ và số cụm $K$. Thuật toán k-means tìm các tâm cụm $y_1,\dots,y_K$ sao cho tổng khoảng cách bình phương từ mỗi điểm tới tâm cụm gần nhất là nhỏ nhất:
\begin{equation}
    E(y_1,\dots,y_K)
    =
    \sum_{i=1}^{m}
    \min_{c\in\{1,\dots,K\}}\|x_i-y_c\|_2^2.
    \label{eq_kmeans_objective}
\end{equation}
Đặt $c_i$ là chỉ số cụm của điểm $x_i$. Khi đó, bài toán có thể viết lại là
\begin{equation}
    E(y_1,\dots,y_K;c_1,\dots,c_m)
    =
    \sum_{i=1}^{m}\|x_i-y_{c_i}\|_2^2.
    \label{eq_kmeans_expanded}
\end{equation}
K-means chính là tối ưu luân phiên giữa hai bước \cite{Lloyd1982}:
\begin{itemize}
    \item \textbf{Bước gán cụm:} cố định các tâm $y_j$, chọn
    \begin{equation}
        c_i=\arg\min_{c\in\{1,\dots,K\}}\|x_i-y_c\|_2^2.
        \label{eq_kmeans_assignment}
    \end{equation}
    \item \textbf{Bước cập nhật tâm:} cố định các $c_i$, cập nhật
    \begin{equation}
        y_j=
        \frac{\sum_{i:c_i=j}x_i}{|\{i:c_i=j\}|}.
        \label{eq_kmeans_centroid}
    \end{equation}
\end{itemize}
Mỗi bước không làm tăng hàm mục tiêu. Tuy nhiên, vì bài toán không lồi, kết quả phụ thuộc vào khởi tạo. Một cách xử lý phổ biến là chạy k-means nhiều lần hoặc dùng k-means++ để chọn tâm ban đầu tốt hơn \cite{Arthur2007KMeansPP}.

\subsection{Lagrangian tăng cường và ADMM}
\label{SUB_SEC_admm}

\subsubsection*{Từ hình phạt bậc hai đến nhân tử Lagrange}
Xét bài toán có ràng buộc đẳng thức
\begin{equation}
    \begin{aligned}
        &\min_x \quad f(x)\\
        &\text{với điều kiện}\quad g(x)=0.
    \end{aligned}
    \label{eq_eq_constrained_problem}
\end{equation}
Một cách xử lý là dùng hình phạt bậc hai:
\begin{equation}
    f_{\rho}(x)=f(x)+\frac{\rho}{2}\|g(x)\|_2^2.
    \label{eq_quadratic_penalty}
\end{equation}
Khi $\rho$ lớn, nghiệm của bài toán phạt có xu hướng thỏa ràng buộc tốt hơn. Nhưng nếu $\rho$ quá lớn, bài toán trở nên kém điều kiện và khó tối ưu.
\medskip

Một cách khác là dùng Lagrangian:
\begin{equation}
    \mathcal{L}(x,\lambda)=f(x)-\lambda^{T}g(x).
    \label{eq_lagrangian_basic}
\end{equation}
Cách này đưa ràng buộc vào thông qua nhân tử Lagrange, nhưng bài toán trở thành bài toán yên ngựa: tối thiểu theo $x$ và tối đa theo $\lambda$.
\begin{cy}
    Hai quy ước dấu $f(x)-\lambda^{T}g(x)$ và $f(x)+\lambda^{T}g(x)$ là tương đương nếu thay $\lambda$ bởi $-\lambda$. Để khớp với phần Lagrangian tăng cường cổ điển trong tài liệu, đoạn này dùng dấu trừ. Riêng công thức ADMM phía sau được viết theo đúng quy ước dấu cộng xuất hiện trong tài liệu.
\end{cy}

\subsubsection*{Phương pháp Lagrangian tăng cường}
Lagrangian tăng cường kết hợp hai ý tưởng trên:
\begin{equation}
    \mathcal{L}_{\rho}(x,\lambda)
    =
    f(x)+\frac{\rho}{2}\|g(x)\|_2^2-\lambda^{T}g(x).
    \label{eq_augmented_lagrangian}
\end{equation}
Thuật toán lặp giữa cập nhật biến chính và cập nhật nhân tử:
\begin{equation}
    \begin{cases}
        x_{k+1}=\arg\min_x \mathcal{L}_{\rho}(x,\lambda_k),\\
        \lambda_{k+1}=\lambda_k-\rho g(x_{k+1}).
    \end{cases}
    \label{eq_augmented_lagrangian_steps}
\end{equation}
Thành phần phạt bậc hai giúp bước tối ưu theo $x$ không đi quá xa khỏi tập ràng buộc, còn nhân tử Lagrange giúp không cần đưa $\rho$ tiến tới vô hạn. Cập nhật nhân tử dùng điểm mới $x_{k+1}$, tức nghiệm vừa thu được từ bước tối ưu theo $x$.

\subsubsection*{ADMM: công thức tổng quát}
ADMM, tức phương pháp nhân tử theo hướng luân phiên, áp dụng ý tưởng Lagrangian tăng cường cho bài toán có hai nhóm biến:
\begin{equation}
    \begin{aligned}
        &\min_{x,z}\quad f(x)+h(z)\\
        &\text{với điều kiện}\quad Ax+Bz=c.
    \end{aligned}
    \label{eq_admm_problem}
\end{equation}
Lagrangian tăng cường tương ứng là
\begin{equation}
    \mathcal{L}_{\rho}(x,z,\lambda)
    =
    f(x)+h(z)
    +\lambda^{T}(Ax+Bz-c)
    +\frac{\rho}{2}\|Ax+Bz-c\|_2^2.
    \label{eq_admm_aug_lag}
\end{equation}
ADMM cập nhật ba bước:
\begin{equation}
    \begin{cases}
        x_{k+1}
        =
        \arg\min_x \mathcal{L}_{\rho}(x,z_k,\lambda_k),\\
        z_{k+1}
        =
        \arg\min_z \mathcal{L}_{\rho}(x_{k+1},z,\lambda_k),\\
        \lambda_{k+1}
        =
        \lambda_k+\rho(Ax_{k+1}+Bz_{k+1}-c).
    \end{cases}
    \label{eq_admm_steps}
\end{equation}
Khác với Lagrangian tăng cường cổ điển, ADMM không tối ưu đồng thời $x$ và $z$. Thay vào đó, thuật toán tối ưu luân phiên từng biến. Việc tách này có thể làm tăng số vòng lặp, nhưng mỗi vòng lặp thường rẻ hơn nhiều và dễ song song hóa hơn \cite{Boyd2011ADMM}.

\subsubsection*{Nhận xét về hội tụ và tham số $\rho$}
Trong trường hợp $f$ và $h$ là các hàm lồi, ràng buộc là tuyến tính và Lagrangian tăng cường có điểm yên ngựa, ADMM có bảo đảm hội tụ tới nghiệm tối ưu toàn cục \cite{Boyd2011ADMM}. Với bài toán không lồi, ADMM vẫn thường được dùng trong thực nghiệm, nhưng lý thuyết hội tụ yếu hơn và kết quả có thể phụ thuộc vào cách tách biến cũng như khởi tạo.
\medskip

Một đặc điểm của ADMM là thuật toán thường giảm nhanh sai số ở giai đoạn đầu, tức nhanh chóng tạo ra nghiệm xấp xỉ tốt, nhưng có thể cần nhiều vòng lặp để đạt độ chính xác rất cao. Vì vậy, trong một số hệ thống, ADMM được dùng để tìm nghiệm ban đầu tốt rồi chuyển sang một phương pháp tinh chỉnh khác nếu cần độ chính xác cuối rất chặt.
\medskip

Tham số $\rho>0$ điều khiển mức phạt đối với vi phạm ràng buộc. Nếu $\rho$ quá nhỏ, biến nguyên thủy có thể tiến chậm tới miền thỏa ràng buộc. Nếu $\rho$ quá lớn, bước cập nhật có thể trở nên kém điều kiện hoặc quá tập trung vào ràng buộc mà giảm ít hàm mục tiêu. Trong triển khai, người ta thường theo dõi phần dư nguyên thủy và phần dư đối ngẫu để điều chỉnh $\rho$: nếu phần dư nguyên thủy lớn hơn nhiều thì tăng $\rho$, còn nếu phần dư đối ngẫu lớn hơn nhiều thì giảm $\rho$.

\subsubsection*{Ví dụ: bình phương tối thiểu không âm bằng ADMM}
Xét bài toán
\begin{equation}
    \begin{aligned}
        &\min_x \quad \|Ax-b\|_2^2\\
        &\text{với điều kiện}\quad x\geq 0.
    \end{aligned}
    \label{eq_nnls_problem}
\end{equation}
Ràng buộc $x\geq 0$ khiến bình phương tối thiểu tuyến tính thông thường không còn áp dụng trực tiếp. Đưa vào biến phụ $z$, bài toán được viết tương đương thành:
\begin{equation}
    \begin{aligned}
        &\min_{x,z}\quad \|Ax-b\|_2^2+h(z)\\
        &\text{với điều kiện}\quad x=z,
    \end{aligned}
    \label{eq_nnls_split}
\end{equation}
trong đó
\begin{equation}
    h(z)=
    \begin{cases}
        0, & z\geq 0,\\
        +\infty, & \text{ngược lại}.
    \end{cases}
    \label{eq_indicator_nonnegative}
\end{equation}
Lagrangian tăng cường là
\begin{equation}
    \mathcal{L}_{\rho}(x,z,\lambda)
    =
    \|Ax-b\|_2^2+h(z)
    +\lambda^{T}(x-z)
    +\frac{\rho}{2}\|x-z\|_2^2.
    \label{eq_nnls_admm_lagrangian}
\end{equation}
Bước cập nhật $x$ thu được bằng cách cho gradient theo $x$ bằng $0$:
\begin{equation}
    (2A^{T}A+\rho I)x_{k+1}
    =
    2A^{T}b+\rho z_k-\lambda_k.
    \label{eq_nnls_x_update_system}
\end{equation}
Do đó,
\begin{equation}
    x_{k+1}
    =
    (2A^{T}A+\rho I)^{-1}(2A^{T}b+\rho z_k-\lambda_k).
    \label{eq_nnls_x_update}
\end{equation}
Bước cập nhật $z$ là phép chiếu lên miền không âm:
\begin{equation}
    z_{k+1}
    =
    \max\left(x_{k+1}+\frac{1}{\rho}\lambda_k,0\right),
    \label{eq_nnls_z_update}
\end{equation}
trong đó phép $\max$ được hiểu theo từng phần tử. Cuối cùng,
\begin{equation}
    \lambda_{k+1}
    =
    \lambda_k+\rho(x_{k+1}-z_{k+1}).
    \label{eq_nnls_lambda_update}
\end{equation}
NNLS lúc này được tách thành ba thao tác rõ ràng: giải hệ tuyến tính, chiếu lên miền không âm và cập nhật biến đối ngẫu.

\subsubsection*{Ví dụ: trung vị hình học bằng ADMM}
Xét lại bài toán trung vị hình học:
\begin{equation}
    \min_x \sum_{i=1}^{N}\|x-x_i\|_2.
    \label{eq_admm_geomedian_original}
\end{equation}
Đặt biến phụ $z_i=x_i-x$, hay tương đương $z_i+x=x_i$, bài toán trở thành
\begin{equation}
    \begin{aligned}
        &\min_{x,\{z_i\}}\quad \sum_{i=1}^{N}\|z_i\|_2\\
        &\text{với điều kiện}\quad z_i+x=x_i,\qquad i=1,\dots,N.
    \end{aligned}
    \label{eq_admm_geomedian_split}
\end{equation}
Với nhân tử $\lambda_i$ cho từng ràng buộc, bước cập nhật $x$ có dạng đóng:
\begin{equation}
    x_{k+1}
    =
    \frac{1}{N}\sum_{i=1}^{N}
    \left(x_i-z_i^{(k)}-\frac{1}{\rho}\lambda_i^{(k)}\right).
    \label{eq_admm_geomedian_x_update}
\end{equation}
Bước cập nhật $z_i$ tách rời theo từng $i$. Đặt
\begin{equation}
    u_i^{(k)}
    =
    x_i-x_{k+1}-\frac{1}{\rho}\lambda_i^{(k)}.
    \label{eq_admm_geomedian_u}
\end{equation}
Khi đó,
\begin{equation}
    z_i^{(k+1)}
    =
    \max\left(1-\frac{1}{\rho\|u_i^{(k)}\|_2},0\right)u_i^{(k)}.
    \label{eq_admm_geomedian_z_update}
\end{equation}
Cập nhật này là phép co vectơ về phía $0$, tương tự soft-thresholding nhưng theo chuẩn Euclid. Cuối cùng,
\begin{equation}
    \lambda_i^{(k+1)}
    =
    \lambda_i^{(k)}+\rho(z_i^{(k+1)}+x_{k+1}-x_i).
    \label{eq_admm_geomedian_lambda_update}
\end{equation}
Ví dụ này cho thấy lợi thế của ADMM: bài toán ban đầu không trơn tại các điểm $x=x_i$, nhưng sau khi tách biến, các bước cập nhật có dạng đơn giản và có thể song song hóa theo từng $i$.
\medskip

Tóm lại, ADMM không chỉ là một thuật toán cụ thể mà còn là một khuôn mẫu thiết kế thuật toán. Khi gặp một bài toán khó, ta tìm cách tách nó thành các khối sao cho mỗi bước cập nhật có thể giải nhanh, có nghiệm đóng hoặc có phép chiếu đơn giản.
